\section{Introduction}
\label{sec:intro}

GPUs and tensor accelerators illustrate how hardware specialization can improve the performance and energy efficiency of data-intensive workloads~\cite{autocomp}. Exploiting these capabilities has traditionally required substantial human effort: experts must organize data movement, select instructions, and tune kernels for the target hardware. Architecture-specific programming interfaces make this work difficult to reuse when moving to another platform. Recent work has begun to automate parts of this process with language-model agents that generate and optimize kernels using hardware knowledge and execution feedback~\cite{kernelbench,autocomp}.

Specialized execution units are also embedded in general-purpose CPUs, where SIMD and vector extensions accelerate data-parallel computation. Their execution models differ substantially: Arm NEON operates on fixed-width vectors, whereas RISC-V Vector (RVV) allows the active vector length to be selected at runtime~\cite{acle,rvv}. A migration between them must account for different loop structures, tail handling, and instruction behavior. RVV is a more recent addition to this ecosystem, and extending its kernel support has motivated work on porting existing vector code~\cite{saltyrn,intrintrans}. C intrinsic interfaces expose these hardware operations within ordinary C programs. Programmers can combine C control flow and pointer operations with vector intrinsics that the compiler lowers to target instructions, avoiding the need to implement an entire kernel in assembly.

This programming interface also gives agents a way to apply their code-generation capabilities to SIMD optimization. Intrinsic descriptions and hardware-specific guidance can inform changes to instruction selection, vector grouping, and loop organization~\cite{saltyrn,autocomp}. However, these changes can alter program behavior on inputs that the available tests do not exercise. Reports of reward hacking in kernel-generation benchmarks further demonstrate that an optimization objective can be satisfied by exploiting weaknesses in its evaluation~\cite{kernelbench-eval}. Existing work addresses correctness through testing and formal validation, including bounded equivalence checking of translated kernels and verification of vectorization in compiler IR~\cite{saltyrn,llm-vectorizer}. Establishing equivalence for the complete calling domain requires reasoning about the kernels' loops and memory accesses while accounting for numerical differences between their ISAs. We seek this guarantee for the actual optimized implementations, including the control flow and intrinsic choices responsible for their performance.

A shared semantic basis for this verification can be constructed at several levels. At the instruction level, ISA specifications describe operations on architectural state, enabling machine-code proofs grounded in detailed hardware semantics~\cite{sail,islaris}. At the intermediate-representation level, vector operations can be expressed in an IR and compared using its computational and memory semantics. LLVM translation validation and verified instruction selection for Hexagon HVX demonstrate uses of this approach~\cite{alive2,rake}. Intrinsic source code provides a third option: interpreting supported intrinsics as C operations turns the comparison into a C program verification problem. The choice determines how vector behavior is connected to the surrounding control flow and memory accesses, and which semantic definitions can be reused across kernels.

Our verification target is functional equivalence under explicit calling contracts: the two kernels must agree on their returned outputs and satisfy their memory specifications for admitted inputs and hardware configurations. For this target, the intrinsic interface exposes the vector operations together with the C loops and pointers that manipulate them. Reasoning from compiled instructions would also require relating architectural states and connecting the machine-code behavior to the source kernels. We instead interpret the intrinsics at the C level, where their numerical and memory effects can be inspected alongside the source program and handled by existing C verification infrastructure. This choice preserves architecture-specific behavior in the intrinsic definitions while giving the complete kernels a common language for verification. The correspondence between those definitions and native instructions remains an explicit assumption of the approach.

We present \tool, a framework for cross-ISA kernel equivalence checking through reusable C semantic shims, with NEON, RVV, AVX, and Hexagon HVX as its target interfaces. \todo{Report the supported intrinsic subsets.} For each kernel pair, \tool binds intrinsic calls to their semantic definitions and states the required output relation in a common specification. QCP symbolically executes the annotated C programs and generates verification conditions for memory accesses, loop invariants, and postconditions, which are discharged through automation and Rocq/Coq proofs~\cite{qcp}. Each implementation is proved against the common specification using its own invariants, so their loops need not advance in lockstep. Implementations, specifications, and proofs remain separate, allowing agents to develop annotations and auxiliary lemmas while retaining the executable kernels. The resulting guarantee is partial correctness relative to the calling contracts and semantic models (Section~\ref{sec:guarantee}).

Human review centers on the calling contract and the evidence produced during verification. A concrete counterexample must first be checked against the original caller's assumptions. If those assumptions exclude the input, the specification can be strengthened to express them, with evidence that all legitimate calls remain covered. If the input is feasible, it identifies a behavioral discrepancy for target-side repair. The revised implementation is then verified again and benchmarked to assess the repair's performance cost. An unsuccessful proof attempt can also reflect a missing invariant or lemma, so unresolved obligations remain distinct from demonstrated discrepancies. Keeping proof changes separate from executable changes makes it possible to review the reason for each repair and its effect on the optimized code.

This equivalence framework supports three uses. First, it checks whether an optimized kernel preserves its reference behavior, within one ISA or across ISAs. Second, it supplies a correctness criterion for porting existing NEON or AVX kernels to RVV or HVX libraries that lack the corresponding operation. Existing implementations then provide both source code for translation and a reference for verification. Third, it enables semantic cross-checking of independently developed implementations of the same operation. As in differential testing, disagreement provides a starting point for investigation; equivalence checking makes the relationship between the implementations an explicit proof obligation~\cite{mckeeman,symdiff}. Examining a concrete discrepancy against the shared contract can identify an implementation defect or an allowed numerical difference, guiding repairs and clarifications to the library's documented behavior.

Our planned evaluation covers these uses through 35 final optimized RVV kernels from SALTyRN, five AVX operations to be ported to RVV and HVX, and corresponding implementations in XNNPACK~\cite{final-kernels}. A composed SIMD program, with softmax as a candidate, will examine whether the kernel specifications support equivalence reasoning across operation boundaries. We will measure verification coverage, proof-development time and token cost, the changes required for verification or repair, and performance against consistent target-side baselines. \todo{Report verified pairs and speedups, completed ports, cross-checking findings and their upstream status, and the composed-program result.} These experiments will assess the extent to which the framework can validate optimized code, reuse semantic definitions across implementations, and support repairs that retain the benefit of optimization.

% Alternative short P8, retained at the author's request; not rendered.
% The framework supports validation of kernel optimizations, verified translation
% across ISAs, and cross-checking of existing implementations. The planned main
% study uses SALTyRN's 35 published final optimized RVV kernels. Additional studies
% target AVX-to-RVV/HVX ports, XNNPACK implementation comparisons, and a composed
% SIMD program. All proof counts, repairs, cost measurements, and speedups remain
% \todo{pending evaluation}; none is asserted as a completed contribution.
